\section{The Guardian Angels}\label{sec:guardians}

God, who governs the whole creation through the ministry of the angels,
gives to every man an angel of his own. That angel keeps him from the
beginning of his life to its end, turns his mind toward the good, holds back
the enemy, carries his prayers to God and joins his own to them, and at the
end of the way brings him home. Scripture gives the doctrine two pillars:
the psalm, ``For he hath given his angels charge over thee; to keep thee in
all thy ways'' (Ps 90[91]:11), and the word of Christ, ``See that you
despise not one of these little ones: for I say to you, that their angels in
heaven always see the face of my Father who is in heaven'' (Matt 18:10). On
these the Fathers built the angel of each soul: the angel of righteousness
of the \emph{Shepherd}, Origen's angels of the little ones, Basil's tutor
and shepherd beside each believer, Jerome's \latin{magna dignitas
animarum}, Chrysostom's ``This is a truth, that each man has an Angel.'' Bernard made
the doctrine a school of reverence and confidence; the Master of the
\emph{Sentences} set it in the schools; Aquinas gave it the eight articles of
q.~113, between the mission of the angels (\S\ref{sec:q112}) and the
assaults of the demons (\S\ref{sec:q114}), where the guardian angel appears
as God's providence exercised toward each immortal person. The Church
confesses the doctrine in her Catechism, keeps the feast of the Guardian
Angels on 2 October, and teaches her children from infancy to pray to their
angel.

\sectionguard
\subsection{The Guardian Angel in Sacred Scripture}\label{sec:gd-scripture}

\paragraph{The angel of the patriarchs.} Abraham sends his servant to find a
wife for Isaac with the assurance that God ``will send his angel before
thee'' (Gen 24:7), and the servant repeats it: ``The Lord \dots\ will send
his angel with thee, and will direct thy way'' (Gen 24:40). Jacob, blessing
the sons of Joseph at the end of his life, joins the God who fed him to the
angel who kept him: ``God that feedeth me from my youth until this day: The
angel that delivereth me from all evils, bless these boys'' (Gen
48:15--16). Chrysostom reads the verse as the proof that ``even from the
beginning, every one of those that were approved had his Angel'' (\emph{Hom.\
in Col.}\ 3).

\paragraph{The angel of the Exodus.} God promises Israel in the desert a
guide and a guard: ``Behold I will send my angel, who shall go before thee,
and keep thee in thy journey, and bring thee into the place that I have
prepared. Take notice of him, and hear his voice, and do not think him one
to be contemned: for he will not forgive when thou hast sinned, and my name
is in him'' (Exod 23:20--21). After the golden calf the promise is renewed, ``my angel shall go before thee'' (Exod 32:34), and
Moses begs for more: ``If thou thyself dost not go before, bring us not out
of this place'' (Exod 33:15), a request Origen and Hilary read as the prayer
of one who seeks the Lord's own keeping beyond an angel's (\emph{Comm.\ in
Matt.}\ XIII.26; \emph{Tract.\ in Ps.}\ 124.6). Exod 23:20 is the \emph{sed
contra} for the mission of the angels (\ST{I q.~112 a.~1}).

\paragraph{The Psalms.} ``The angel of the Lord shall encamp round about them
that fear him: and shall deliver them'' (Ps 33[34]:8). Psalm 90 promises the
man who dwells in the aid of the Most High a shield against ``the terror of
the night'' and ``the noonday devil'' (vv.~5--6), and gives the reason: ``For he hath given his angels
charge over thee; to keep thee in all thy ways. In their hands they shall
bear thee up: lest thou dash thy foot against a stone. Thou shalt walk upon
the asp and the basilisk: and thou shalt trample under foot the lion and the
dragon'' (vv.~11--13). The first verse is the \emph{sed contra} of
Aquinas's first article (\ST{I q.~113 a.~1}); another psalm, that God himself
``shall neither slumber nor sleep, that keepeth Israel'' (Ps 120[121]:4),
supplies the article's second objection.

\paragraph{Tobias.} The Book of Tobias shows the guardian's office at length.
The young Tobias, sent on a long road, finds ``a beautiful young man,
standing girded, and as it were ready to walk,'' ``not knowing that he was
an angel of God'' (Tob 5:5--6), and his father comforts the weeping mother:
``I believe that the good angel of God doth accompany him, and doth order all
things well that are done about him'' (Tob 5:27). Raphael delivers the boy
from the fish (Tob 6:2--4), binds the demon that had slain Sara's husbands
(Tob 8:3), and brings him home; then he reveals the hidden side of his care:
``When thou didst pray with tears, and didst bury the dead \dots\ I offered
thy prayer to the Lord,'' and ``when I was with you, I was there by the will
of God'' (Tob 12:12, 18). John Paul~II reads the book as \latin{così
significativa circa l'affidamento agli angeli dei piccoli figli di Dio,
sempre bisognosi di custodia, di cura e di protezione}: so significant for
the entrusting to the angels of God's little children, always in need of
keeping, care, and protection (general audience, 6 August 1986, n.~8).
Judith returns from the camp of Holofernes confessing the same keeping: ``his
angel hath been my keeper both going hence, and abiding there, and returning
from thence hither'' (Jdt 13:20).

\paragraph{Daniel.} The tenth chapter of Daniel shows the angels over nations.
The angel sent to Daniel says, ``from the first day that thou didst set thy
heart to understand \dots\ thy words have been heard: and I am come for thy
words. But the prince of the kingdom of the Persians resisted me one and
twenty days: and behold Michael, one of the chief princes, came to help me''
(Dan 10:12--13), and ``none is my helper in all these things, but Michael
your prince'' (10:21). Michael is ``the great prince, who standeth for the
children of thy people'' (Dan 12:1). The Septuagint of the Song of Moses,
which Basil, Hilary, and Jerome read, says that the Most High set the
bounds of the nations according to the number of the angels of God (Deut
32:8); the Vulgate has ``the children of Israel.'' The chapter is the
\emph{sed contra} of Aquinas's eighth article (\ST{I q.~113 a.~8}).

\paragraph{The Gospels and Acts.} At the pinnacle of the Temple the devil
quotes the psalm of the guardians to Christ (Matt 4:6); when the tempter
departs, ``angels came and ministered to him'' (4:11), and in the agony ``there
appeared to him an angel from heaven, strengthening him'' (Luke 22:43).
Christ's saying on the little ones (Matt 18:10) is the foundation of the whole
doctrine: the little ones have ``their angels,'' and those angels at the same
time ``always see the face'' of the Father. There is ``joy before the
angels of God upon one sinner doing penance'' (Luke 15:10), and the beggar
Lazarus ``was carried by the angels into Abraham's bosom'' (Luke 16:22). When
an angel has led Peter out of Herod's prison (Acts 12:7--11), the disciples in
the house of Mary hear his voice at the gate and say, ``It is his angel''
(12:15); on the ship to Rome Paul speaks of ``an angel of God, whose I am and
whom I serve'' (Acts 27:23). The Epistle to the Hebrews states the general
law: ``Are they not all ministering spirits, sent to minister for them who
shall receive the inheritance of salvation?'' (Heb 1:14).

\sectionguard
\subsection{The Fathers and the Angel of Each Soul}\label{sec:gd-fathers}

The Fathers read these texts together and drew from them the angel of each
soul. Each Father's fuller teaching on the angels is expounded in its own
place (\S\S\ref{sec:fathers-before}, \ref{sec:fathers-latin},
\ref{sec:fathers-after}). Together they teach that an angel is given to each
soul, speaks within it through its thoughts, prays with it and for it, keeps
it on the road, and brings it to the judgment.

\subsubsection{The Shepherd, Clement, and Origen}

The \emph{Shepherd} of Hermas gives the Church her earliest teaching that an
angel attends each man and is known by the thoughts he raises: ``There are
two angels with a man---one of righteousness, and the other of iniquity.''
When the angel of righteousness ``ascends into your heart, forthwith he
talks to you of righteousness, purity, chastity, contentment, and of every
righteous deed and glorious virtue''; the other is known by anger,
bitterness, and craving, and the rule is, ``Now that you know his works,
depart from him'' (\emph{Mand.}\ VI.2). The good angel speaks within; the man
discerns and consents. Cassian carries the two angels to the monks, and the
Master of the \emph{Sentences} to the schools (\S\ref{sec:gd-masters}).
Clement of Alexandria sets the angel of the individual within the angels of
the nations: ``For regiments of angels are distributed over the nations and
cities. And, perchance, some are assigned to individuals'' (\emph{Strom.}\
VI.17).

Origen's commentary on Matt 18:10, expounded with his other teaching in
\S\ref{sec:fathers-before}, likens the guardians to the nursing-fathers and
tutors of the imperfect, whom the perfect outgrow to be ``governed by the
Lord Himself'' (\emph{Comm.\ in Matt.}\ XIII.26), and raises the question
whether the angels take up their charge ``from what time `by the laver of
regeneration' \dots\ or, whether from birth'' (XIII.27), which Aquinas
decides with Jerome (a.~5). The homilies on Numbers carry the guardianship to
the harvest and the judgment. The Lord's field is the world, \latin{ager
autem non terrae solum, sed corda intelliguntur humana, quem agrum angeli Dei
susceperint excolendum}: not the earth only but human hearts, a field the
angels of God have taken in hand to till. They keep as their own fruit those
still under tutors, and offer the hearts brought to perfection as
first-fruits to the great High Priest, as Cornelius was offered by his angel
(\emph{Hom.\ in Num.}\ XI.3; Acts 10:4). At the end they come with their
harvest: \latin{Igitur unusquisque angelorum in consummatione saeculi aderit
in iudicio, perducens secum eos, quibus praefuit, quos adiuvit, quos
instruxit, pro quibus semper vidit faciem patris, qui est in coelis}. Each
angel will stand at the judgment at the end of the age, bringing those over
whom he presided, whom he helped and taught, for whom he always saw the
Father's face. \latin{Et puto etiam ibi inquisitionem futuram, utrum culturae
hominum angelus defuerit, an culturae angelicae nequaquam digne segnitia
humana responderit}: there will be an inquiry whether the angel failed in
tilling men or human sloth answered unworthily to the angel's tilling
(XI.4). Aquinas meets the sentence in the Gloss on Numbers and answers it in
his article on the angels' grief (a.~7 obj.~4). The \emph{De
principiis} proposes, among readings of the unfaithful stewards who ``are to
be divided'' (Matt 24:51), that since ``each believer, although the humblest in the Church, is
said to be attended by an angel,'' the angel is withdrawn at the end from the
one who ``has not faithfully observed the admonitions of the angel allotted it
by God'' (II.10.7). Origen's pupil Gregory the Wonderworker gave thanks for
``that holy angel of God who fed me from my youth,'' who carried him past
Berytus, where he meant to study the laws, and brought him to Origen
(\emph{Pan.\ Or.}\ 4--5).

\subsubsection{Basil, Gregory of Nyssa, and Chrysostom}

Basil gives the sentence the Catechism makes its own. The angels are one in
nature and differ in office: some preside over nations, some accompany each
of the faithful, and the angel of a nation has the greater dignity; and
\gk{to de syneinai hekast\=o t\=on pist\=on angelon, hoion paidag\=ogon
tina kai nomea t\=en z\=o\=en dieuthynonta, oudeis anterei}, that an angel
is with each of the faithful, as a kind of tutor and shepherd directing his
life, no one will gainsay who remembers the Lord's word on the little ones
(\emph{Adversus Eunomium} III.1; PG~29, 656B--657A). The \gk{paidag\=ogos}
leads a child to school; the \gk{nomeus} herds the flock. The Catechism
renders the sentence, ``Beside each believer stands an angel as protector and
shepherd leading him to life'' (CCC~336). In the homily on Psalm 33 Basil
warns that the guardian can be driven off: \gk{panti pepisteukoti eis ton
Kyrion angelos paredreuei, ean m\=epote auton h\=emeis ek t\=on pon\=er\=on
erg\=on apodi\=ox\=omen}, an angel sits beside everyone who has believed in
the Lord, unless we drive him away by evil works, as smoke drives off bees;
the one angel is likened to a whole camp because he walls us round on every
side, and so ``a thousand shall fall at thy side'' (\emph{Hom.\ in Ps.}\ 33.5;
Ps 90[91]:7, 11). Basil's teaching is set out at length in
\S\ref{sec:fb-basil}; Aquinas's sixth article answers the question his image
raises.

Gregory of Nyssa hands on the guardian as a tradition of the fathers. Reading
Aaron's coming out to meet Moses (Exod 4:27), he writes, \gk{Logos tis estin
ek patrik\=es paradose\=os to piston ech\=on}, there is a teaching,
trustworthy from the tradition of the fathers, that when our nature fell
into sin God did not leave our fall without providence, but set beside each
one's life an incorporeal angel \gk{eis symmachian}, as an ally, while the
corrupter set an evil demon against it; man stands between the two and gives
the victory by his choice. The ally was given at our first birth and shows
himself when we set out for the higher life, for \gk{adelphos gar tropon
tina kata to logikon te kai noeron t\=es psych\=es tou anthr\=opou, ho
angelos}, the angel is in a certain way a brother to the rational and
intellectual part of man's soul (\emph{De vita Moysis} II; PG~44, 337--340;
\S\ref{sec:fb-nyssa}).

Chrysostom, whose teaching is set out in \S\ref{sec:fb-chrysostom}, states
the doctrine on Acts 12:15 without qualification: ``This is a truth, that
each man has an Angel'' (\emph{Hom.\ in Act.}\ 26). On Matt 18:10 he draws
the conclusion ``that the saints have angels, or even all men,'' and adds that
the Lord speaks ``of angels that are greater than others'' (\emph{Hom.\ in
Matt.}\ 59.4), the first objection of Aquinas's third article. On Colossians
he gives the guardian a place in the order of reconciliation: ``At first the
Angels were according to the number of the nations; but now, not to the
number of the nations, but that of the believers. \dots\ If then we have
Angels, let us be sober, as though we were in the presence of tutors; for
there is a demon present also. Therefore we pray, asking for the Angel of
peace'' (\emph{Hom.\ in Col.}\ 3).

\subsubsection{The Latin Fathers}

Hilary teaches that our weakness needed the angels' guard: \latin{neque enim
infirmitas nostra nisi datis ad custodiam angelis tot tantisque spiritalium
caelestium nequitiis obsisteret}, our weakness could not have withstood so
many spiritual wickednesses unless angels had been given for its keeping
(\emph{Tract.\ in Ps.}\ 134.17); yet with Moses he prefers the Lord's keeping,
\latin{bonum quidem praesidium angeli, sed melius domini} (124.6;
\S\ref{sec:fl-hilary}). Ambrose draws a duty of prayer: ``The angels must be
entreated for us, who have been to us as guards'' (\emph{De viduis} 9.55;
\S\ref{sec:fl-ambrose}). Jerome gives the schools their authority:
\latin{Magna dignitas animarum, ut unaquaeque habeat ab ortu nativitatis in
custodiam sui angelum delegatum} (\emph{In Matth.}\ III, on 18:10; PL~26).
Great is the dignity of souls, that each should have from the beginning of
its birth an angel delegated to its keeping. The sentence is the \emph{sed
contra} of three of Aquinas's articles (aa.~2, 4, 5; \S\ref{sec:fl-jerome}).
Abbot Serenus in Cassian gathers the Scripture of the guardian into one
chapter: ``Holy Scripture bears witness that two angels, a good and a bad one,
cling to each one of us,'' citing Matt 18:10, Ps 33[34]:8, Acts 12:15, and
``the book of the Shepherd'' (\emph{Conl.}\ VIII.17; \S\ref{sec:fl-cassian}).

Gregory the Great gives the lowest order its office, \latin{Hi autem qui
minima nuntiant, angeli \dots\ vocantur}, those who announce the least
things are called angels; the Powers restrain the hostile spirits \latin{ne
corda hominum tantum tentare praevaleant quantum volunt}, lest they tempt the
hearts of men as much as they would (\emph{Hom.\ in Ev.}\ 34.8, 10). His
\emph{Moralia} supplies the angels of the nations and the reading of
Daniel's prince of Persia that Aquinas follows (XVII.12.17;
\S\ref{sec:fa-gregory}), and states the rule of every trial God permits:
``the dispensation of God both while guarding, forsakes his elect servant,
and while forsaking, guards him'' (\emph{Moralia} III, \S6, on Job 2:6).
Isidore handed on both teachings, that each nation has its presiding angels
and that all men have angels, by Matt 18:10 and Acts 12:15 (\emph{Sent.}\
I.10.20--21; \S\ref{sec:fa-isidore}).

\sectionguard
\subsection{Bernard and the Masters of the Twelfth Century}\label{sec:gd-masters}

\paragraph{Bernard on the psalm \emph{Qui habitat}.} Because the verses of
Psalm 90 were in the mouths of the faithful above all in Lent, Bernard
preached seventeen Lenten sermons on it to his monks at Clairvaux. The
eleventh and twelfth, on ``he hath given his angels charge over thee,'' are
the Church's classic meditation on the guardian angel.\footnote{Bernard's
Latin follows Migne, PL~183, cols.~225--237.} The eleventh begins from the
precision of the verse: \latin{Nam neque angelis mandavit ut in omnibus viis
custodiant nos, sed in omnibus viis nostris}; he has not charged his angels
to guard us in all ways, but in all our ways (\emph{In Ps.\ Qui habitat} 11.2).
The ways of men are need and desire, the ways of the demons presumption and
obstinacy; the ways of the holy angels are those the Only-begotten showed
when he said that the angels would ascend and descend upon the Son of man
(John 1:51):
\begin{quote}
\latin{Sic beati illi spiritus ascendunt per contemplationem Dei, descendunt
per compassionem tui, ut custodiant te in omnibus viis tuis. Ascendunt ad
vultum ejus, descendunt ad nutum ejus; quoniam angelis suis mandavit de te.
Nec tamen vel descendendo visione gloriae fraudantur, quia semper vident
faciem Patris.} (11.6)
\end{quote}
\noindent The blessed spirits ascend by contemplation of God and descend by
compassion for you, to keep you in all your ways; they ascend to his face,
they descend at his nod; and in descending they are not deprived of the
vision of glory, for they always see the face of the Father. Ascending they
seek the truth; \latin{Cum vero descendunt, faciunt nobiscum misericordiam,
ut custodiant nos in omnibus viis nostris}, descending they do mercy with us,
to keep us in all our ways (11.10). The
angel's compassion is his condescension in mercy, joined to the unbroken
vision of the Father's face; Gregory had given the same doctrine of the
angels who are sent and yet assist (\S\ref{sec:q112}).

The twelfth sermon turns from the angels' ways to the man they keep. God has
sent his Son and his Spirit and promised his face, and lest anything in
heaven be idle in our regard, \latin{beatos illos spiritus propter nos mittis
in ministerium, custodiae nostrae deputas, nostros jubes fieri paedagogos}:
you send those blessed spirits to minister for us, depute them to our
guard, bid them become our tutors; \latin{illos utique spiritus tam felices,
et tuos ad nos, et nostros ad te angelos facis}, you make those happy
spirits your angels to us and our angels to you (12.3); \latin{Mira
dignatio, et vere magna dilectio charitatis!} (12.4). Then he draws the
practical conclusion:
\begin{quote}
\latin{Quantam tibi debet hoc verbum inferre reverentiam, afferre
devotionem, conferre fiduciam! reverentiam pro praesentia, devotionem pro
benevolentia, fiduciam pro custodia. \dots\ In quovis diversorio, in quovis
angulo, angelo tuo reverentiam habe.} (12.6)
\end{quote}
\noindent The word should bring you reverence for their presence, devotion
for their goodwill, confidence for their guardianship; in every lodging, in
every corner, have reverence for your angel. Would you dare in his presence
what you would not dare in mine? Faith proves what the eye does not see:
\latin{Adsunt igitur, et adsunt tibi, non modo tecum, sed etiam pro te.
Adsunt ut protegant, adsunt ut prosint} (12.6); they are present, present to
you, not only with you but for you, present to protect and to profit you.
The honour belongs to God, who gave the command, yet gratitude is owed to
those who obey it: \latin{Simus ergo devoti, simus grati tantis custodibus: redamemus eos,
honoremus eos quantum possumus, quantum debemus} (12.7). They are our future
co-heirs, \latin{interim vero actores et tutores a Patre positos, et
praepositos nobis}, meanwhile stewards and tutors set over us by the Father
(Gal 4:1--2).
\latin{Fideles sunt, prudentes sunt, potentes sunt: quid trepidamus? Tantum
sequamur eos, adhaereamus eis} (12.8): they are faithful, prudent, and
mighty; why do we tremble? They \latin{deducent parvulum, qua potest
parvulus ambulare}, leading the little one where a little one can walk
(12.8). When temptation presses, \latin{invoca custodem
tuum, ductorem tuum}; call on your guardian and your guide. He does not
sleep, \latin{etsi ad tempus quandoque dissimulet} (12.9), though at times
he seems to look away, lest you fall more dangerously by not knowing whose
hands held you. The sermon ends with the counsel: \latin{Habetote familiares angelos, fratres
mei, frequentate eos sedula cogitatione et devota oratione, qui semper vobis
adsunt ad custodiam et consolationem} (12.10). Make the angels your
familiars, brethren; visit them often by earnest thought and devout prayer,
for they are always with you to guard and console. The thirteenth sermon
reads ``in their hands they shall bear thee up'' of the end of the road: the
angels keep us in our ways, and when the way, which is life, is finished,
they take us up in their hands, as they carried Lazarus (13.1; Luke 16:22).

In \emph{On Consideration} Bernard supposes that ``they are called Angels who
are believed to have been given as guardians of individual men''
(\emph{De consideratione} V.4.8), a sentence Aquinas makes a \emph{sed
contra} for the guardianship of the lowest order (\emph{In II Sent.}\ d.~11
q.~1 a.~2 s.c.~2). The saints who took up Bernard's devotion are treated in
\S\ref{sec:saints}.

\paragraph{Honorius and the \emph{Elucidarium}.} The dialogue of master and
disciple that Suárez still cites as Anselm's, and that Migne prints among the
works of Honorius Augustodunensis, has a chapter \emph{De angelis
custodibus}. The
disciple asks whether men have guardian angels, and the master answers:
\latin{Unicuique genti, unicuique civitati praesunt Angeli, qui jura, leges,
mores juste dispensant et ordinant. Unaquaeque etiam anima, dum in corpus
mittitur, angelo committitur, qui eam semper ad bonum incitet, et omnia
opera ejus Deo et angelis in coelis referat} (\emph{Elucidarium} II.28; PL~172,
1154--1155). Angels preside over every nation and city, and every soul,
when it is sent into the body, is committed to an angel who always spurs it
to the good and reports its works in heaven; since God knows all things, the
reporting is the angels' rejoicing over our progress (Luke 15:10). Asked
whether the angels
are always on earth with those they guard, the master answers,
\latin{Cum opus fuerit, in auxilium veniunt, maxime cum precibus fuerint
invitati}, they come to help when there is need, above all when invited by
prayers; they need no time to come, and \latin{gloria intima non
fraudantur, quia semper vident faciem Patris, quocunque mittantur}, they
lose nothing of their inner glory, for wherever they are sent they always see
the Father's face (II.28). The next chapter adds that a demon overcome in
combat by a just man \latin{mox ab angelo custode ejus in abyssum
retruditur}, is thrust at once into the abyss by the just man's guardian
angel (II.29).

\paragraph{The Master of the \emph{Sentences}.} Peter Lombard gives the
doctrine its place in the schools in the eleventh distinction of his second
book, under the title \latin{Quod singulae animae habeant Angelum bonum ad
custodiam, malum ad exercitium}: that each soul has a good angel for its
guard and an evil one for its exercise. \latin{Illud quoque sciendum est,
quod Angeli boni deputati sunt ad custodiam hominum, ita ut quisque electorum
habeat Angelum, ad sui profectum atque custodiam specialiter delegatum}: the
good angels are deputed to the guardianship of men, so that each of the
elect has an angel specially delegated for his progress and guard. He quotes
Jerome, cites Gregory for the good angel and the bad, and names the angels of
Tobias and of Peter (\emph{Sent.}\ II d.~11 c.~1). Reckoning the elect equal
in number to the good angels and all men more numerous, he answers that one
angel is deputed to many, at the same time or in succession:
\latin{Nec est mirandum unum Angelum pluribus hominibus ad custodiam deputari,
cum uni homini plurium custodia deputetur} (\emph{ibid.}); it is no wonder
that one angel should guard many men, when one man is given the care of
many. Aquinas's commentary on this distinction, written more than a decade
before the \emph{Prima pars}, is the first form of q.~113, and his
\emph{Summa} answers the Master on the number (a.~2). On the evil angel
\latin{ad exercitium} the commentary adds that the fight is just, for a free will
cannot be compelled to sin: \latin{debilis est hostis qui non potest vincere
nisi volentem; et super hoc homini adest praesidium Angeli, et auxilium
divinum, si suscipere velit} (\emph{In II Sent.}\ d.~11 q.~1, expositio
textus); he is a weak enemy who can conquer only the willing, and besides,
man has the protection of his angel and the help of God, if he will accept
it.

\sectionguard
\subsection{The Church's Confession and Prayer}\label{sec:gd-church}

The Catechism of the Catholic Church, after the angels' part in the life of
Christ and of the Church and the Church's invocation of them in the Roman
Canon and the funeral rite, teaches: ``From infancy to death human life is
surrounded by their watchful care and intercession. `Beside each believer
stands an angel as protector and shepherd leading him to life.' Already here
on earth the Christian life shares by faith in the blessed company of angels
and men united in God'' (CCC~336, citing Matt 18:10 and Basil,
\emph{Adv.\ Eun.}\ III.1).

John Paul~II drew the doctrine from Heb 1:14, Ps 90:11--12, Matt 18:10, and
the angels' care for the apostles in Acts, and concluded: \latin{Perciò la Chiesa confessa la sua fede
negli angeli custodi, venerandoli nella liturgia con una festa apposita, e
raccomandando il ricorso alla loro protezione con una preghiera frequente,
come nell'invocazione dell'``Angelo di Dio''} (general audience, 6 August
1986, n.~7). Therefore the Church confesses her faith in the guardian angels,
venerating them in the liturgy with a feast of their own and recommending
recourse to their protection by frequent prayer, as in the \latin{Angele
Dei}, a prayer that, he adds, seems to treasure Basil's words.

The Directory on Popular Piety and the Liturgy (Congregation for Divine
Worship, 2001) lists the days on which the Church keeps the angels, ``29
September (feast of the Archangels Michael, Gabriel and Raphael), 2 October
(the Guardian Angels)'' (n.~215), and describes the devotion that flows from
the doctrine as ``a certain form of the Christian life'' marked by gratitude
to God for these spirits placed at man's service, by devotion from ``living
constantly in the presence of the Holy Angels of God,'' and by ``serenity
and confidence in facing difficult situations''; among the prayers to the
guardian angels ``the Angele Dei is especially popular, and is often recited
by families at morning and evening prayers'' (n.~216), and it names Basil
and Bernard as the devotion's masters. The Directory guards the devotion against
reading the events of life so ``as to ascribe all setbacks to the Devil and
all success to the Guardian Angels,'' and discourages assigning names to the
angels except Gabriel, Raphael, and Michael (n.~217). The feast of 2 October
and its history in the Roman calendars are treated in \S\ref{sec:liturgy};
the \latin{Angele Dei}, its text, and the regulation of devotion to the
angels in \S\ref{sec:devotion}.

\sectionguard
\subsection{The Guarding of the Good Angels (\ST{I q.~113})}\label{sec:q113}

After the angels' action on bodies and on men and their mission, Aquinas
turns to \latin{custodia bonorum Angelorum}, the guardianship of the good
angels, and then to the assault of the bad (q.~114). The eight articles ask
whether men are guarded by angels; whether each man has his own angel;
whether guardianship belongs only to the lowest order; whether every man has
an angel guardian; when the guardianship begins; whether the angel always
guards his charge; whether he grieves at the loss of the one guarded; and
whether there is strife among the angels by reason of their guardianship
(\ST{I q.~113 pr.}). The commentary on the \emph{Sentences} had asked five
questions on the same distinction of the Lombard and a sixth, on the fight
among the angels, in the next question (\emph{In II Sent.}\ d.~11 q.~1
aa.~1--5; q.~2 a.~5); the \emph{Summa} adds the guardian of each man, and
separates the beginning of the guardianship from the question of who
receives it.

\paragraph{Men are guarded by angels (a.~1).} The objections argue that a
guardian is given to those who cannot guard themselves, as children and the
sick, whereas man guards himself by free will and knows by natural law; that
a strong guard makes a weak one superfluous, and ``He shall neither slumber
nor sleep, that keepeth Israel'' (Ps 120[121]:4); and that so many perish
daily whom the angels could help by appearing or by miracles that, if men
were in their keeping, the angels would be negligent. The \emph{sed contra}
is the psalm: ``He hath given His angels charge over thee, to keep thee in
all thy ways'' (Ps 90[91]:11). Aquinas answers from the order of providence:
``in all things the movable and variable are moved and regulated by the
immovable and invariable,'' bodies by spiritual substances, the lower bodies
by the higher, and our own conclusions by principles we hold unchangeably.
``It is moreover manifest that as regards things to be done human knowledge
and affection can vary and fail from good in many ways; and so it was
necessary that angels should be deputed for the guardianship of men, in
order to regulate them and move them to good'' (\latin{per quos regularentur
et moverentur ad bonum}). Free will avoids evil ``to a certain degree, but
not in any sufficient degree,'' for the affections are weak toward the good
through the passions, and natural knowledge of the law fails in its
application to particular deeds: ``The thoughts of mortal men are fearful,
and our counsels uncertain'' (ad~1; Wis 9:14). God guards man in two ways.
He inclines the affection to the good immediately, ``by infusing into him
grace and virtues''; he guides reason in finding the ways to the good ``as
his universal instructor, Whose precepts reach man by the medium of the
angels'' (ad~2; \S\ref{sec:q111}). Men perish by their own malice, not by
the angels' negligence; they depart from the good angels' prompting, which
is given invisibly ``when they enlighten man that he may do what is right,''
as they depart from the natural instinct of the good, and the angels'
visible appearances come ``from a special grace of God,'' like miracles
(ad~3).

The \emph{Sentences} commentary adds the reasons of fittingness. God is the
first and principal guardian, and he executes his providence through the
angels not from any insufficiency but \latin{propter ordinem suae
sapientiae}, for the order of his wisdom: it befits the angels
\latin{ut scilicet eis haec dignitas non negetur, quod sint duces hominum
reductionis in Deum}, that this dignity be not denied them, to be the
leaders of men's return to God; and it befits men, who in the state of
imperfection are like children and so receive guardians \latin{tamquam
pueris} (\emph{In II Sent.}\ d.~11 q.~1 a.~1 ad~1; 1 Cor 13:10--11). God
alone gives grace and glory, as the artist gives the knife its form; the
angels prepare men for grace by illuminations, as fire and hammer dispose the
iron (ad~2). They carry our prayers to God
not to teach him but to pray for us, as the priest offers the people's
prayer (ad~3). They do not compel virtue, \latin{quia sic periret libertas
arbitrii et ratio meriti}, for freedom and merit would perish (ad~5), and
they seldom appear, because a vision above nature overwhelms the mind into
assent and takes from man the good of his own inquiry; ordinarily they
instruct us in our own mode, \latin{illustrando phantasmata, et confortando
lumen intellectus nostri, et excitando ad aliquid rectius considerandum},
enlightening the images of the imagination, strengthening the light of our
understanding, and stirring us to consider something more rightly (ad~6).
The \emph{De veritate} argues from the guardianship that the angels know
singulars, \latin{singulares hominum actus, et alia huiusmodi quae spectant ad
officium custodiae}, the singular acts of men and whatever belongs to the
office of guardianship (\emph{De ver.}\ q.~8 a.~11; \ST{I q.~57 a.~2},
\S\ref{sec:q57}). Suárez gathers the offices of the
guardian from Scripture and the Fathers: to turn aside dangers of body and
soul, to stir the soul to good and warn it of evil, to restrain the demons,
to present our prayers to God, to pray for his ward, and when providence
requires, to correct and chastise him medicinally (\emph{De angelis}
VI.19.1--7).

\paragraph{Each man has his own angel (a.~2).} The objections argue that one
man can guard many and an angel is stronger than a man; that since the
angels are unequal there is only one angel with none between him and men,
who alone can guard men immediately, for Dionysius says the lower are
brought to God through the higher; and that the greater angels have greater
offices, whereas men are equal by nature. The \emph{sed contra} is Jerome:
``Great is the dignity of souls, for each one to have an angel deputed to
guard it from its birth.'' Aquinas determines, \latin{singulis hominibus
singuli Angeli ad custodiam deputantur}: ``Each man has an angel guardian
appointed to him.'' The reason lies in the object of providence. Men are
incorruptible not only in their species but in the proper form of each, the
rational soul, ``which cannot be said of other incorruptible things.''
Providence is chiefly concerned with what abides for ever, and orders what
passes to what abides; so providence stands toward each man as it stands
toward each genus or species of corruptible things. Gregory assigns the
orders to different genera of things, the Powers to restrain the demons, the
Virtues to work miracles in bodies, and it is probable that different
angels of the same order preside over different species; ``Hence it is also
reasonable to suppose that different angels are appointed to the
guardianship of different men.'' The first reply distinguishes two
guardianships. A man may be guarded as an individual, and then ``to one man
one guardian is due; and sometimes several are appointed to guard one''; or
as part of a community, and then one man may guard many in what concerns
their external works, by which others are edified or scandalized. But the
angels guard men ``also as regards invisible and occult things, concerning
the salvation of each one in his own regard'' (\latin{quantum ad invisibilia
et occulta, quae pertinent ad singulorum salutem secundum seipsos}); hence
each has his own (ad~1). The chain of illumination does not forbid it: an
angel may enlighten a man immediately and yet have angels beneath him whom
he enlightens (ad~2; \S\ref{sec:q106}). And though men are equal in nature,
providence orders some to greater things and some to lesser, ``With much
knowledge the Lord hath divided them, and diversified their ways'' (ad~3;
Ecclus 33:11--12); it is a greater office to guard one man than another.

The Lombard had allowed one angel to guard many men at once; the first
reply meets his argument from the good that one man does for many. Suárez,
for whom the guardianship of wayfaring men is \latin{de fide certum} and the
angel of each man a Catholic assertion received by the whole Church
(\emph{De angelis} VI.17.6, 8), judges more likely the opinion \latin{unum
Angelum nunquam custodire simul, nisi unum hominem}, that an angel never
guards more than one man at a time, since an angel cannot be
in two places at once to help two men far apart and the greater number shows
more of God's liberality; that one angel guards many men in succession he
thinks probable (VI.17.10--11). The one guardian does not leave a man to
one angel alone: many angels led Lot out of Sodom, Michael came to help
Gabriel, and, in the words of Origen that Suárez quotes against Calvin,
\latin{Nunquam solus est justus \dots\ sed exercitus ei virtutum praesto
est}, the just man is never alone, but an army of powers stands by him
(VI.17.9).

\paragraph{The guarding of single men belongs to the lowest order (a.~3).}
The first objection is Chrysostom's: the angels of Matt 18:10 are ``not of
any angels but of the highest.'' The second argues from Heb 1:14 that the
mission of the angels is ordered to the guarding of men, and five orders are
sent (\ST{I q.~112 a.~4}); the third, that guarding men requires above all
restraining the demons and working miracles, which Gregory gives to the
Powers and the Virtues. The \emph{sed contra}: the psalm attributes the
guarding of men to ``angels,'' who are the lowest order ``according to
Dionysius'' (\emph{CH}~5 and~9; \S\ref{sec:ch9}). Aquinas answers with the
distinction of a.~2. Particular guardianship, by which one angel keeps one
man, ``belongs to the lowest order of the angels, whose place it is,
according to Gregory, to announce the `lesser things'; for it seems to be
the least of the angelic offices to procure what concerns the salvation of
only one man.'' Universal guardianship is multiplied according to the
orders, ``For the more universal an agent is, the higher it is'': the
guardianship of the human race belongs to the Principalities ``or perhaps
to the `Archangels,' whom we call the angel princes. Hence, Michael, whom we
call an archangel, is also styled `one of the princes' (Daniel 10:13).'' The
Virtues guard all corporeal creatures, the Powers guard over the demons, and
the Principalities over the good spirits, ``according to Gregory's
opinion.'' Chrysostom ``can be taken to mean the highest in the lowest order
of angels; for, as Dionysius says \dots\ in each order there are first,
middle, and last. It is, however, probable that the greater angels are
deputed to keep those chosen by God for the higher degree of glory''
(\latin{Est autem probabile quod maiores Angeli deputentur ad custodiam
eorum qui sunt ad maiorem gradum gloriae a Deo electi}; ad~1; \emph{CH}
10.2). Not all who are sent guard single men; some orders have a universal
guardianship (ad~2). The lower angels exercise the offices of the higher by
sharing their gifts, ``and they are executors of the superiors' power; and
in this way all the angels of the lowest order can coerce the demons, and
work miracles'' (ad~3).

The \emph{Contra gentiles} gives the three lowest orders their objects in
human things. The Principalities have the common good of a city or nation,
so that \latin{dispositio regnorum, et mutatio dominationis a gente in
gentem, ad ministerium huius ordinis pertinere oportet}; the Archangels a
good that belongs to one and profits many, as the things of faith and
divine worship, which Gabriel announced to the Virgin; and the Angels the
good that belongs to each singly: \latin{Et huiusmodi ad ordinem pertinent
Angelorum, de quibus Gregorius dicit quod infima nuntiant: unde et hominum
custodes esse dicuntur} (\emph{SCG} III.80). Dionysius himself had said that
the order of Principalities, Archangels, and Angels ``presides, through each
other, over the Hierarchies amongst men,'' and that ``the Word of God has
assigned our Hierarchy to Angels, by naming Michael as Ruler of the Jewish
people, and others over other nations'' (\emph{CH} 9.2). The
\emph{Sentences} commentary calls the procuring of what concerns determinate
persons \latin{ultimum et minimum in actibus Angelorum}, the last and least
of the angels' acts (\emph{In II Sent.}\ d.~11 q.~1 a.~2). Suárez repeats that the guardians are not from the four
highest orders, and that particular guardianship is committed to the lowest
(\emph{De angelis} VI.18.3, 6), as Basil had ranked the angel of a nation
above the angel of one man (\emph{Adv.\ Eun.}\ III.1).

\paragraph{Every man has a guardian while he is a wayfarer (a.~4).} The
objections exclude Christ, made ``in the likeness of men'' (Phil 2:7) yet
greater than the angels; Adam in innocence; and the foreknown to damnation,
the infidels, and the Antichrist (2 Thess 2:9). The \emph{sed contra} is Jerome again, that ``each
soul has an angel appointed to guard it.'' Aquinas answers with the image
that gives the doctrine its shape: ``Man while in this state of life, is, as
it were, on a road by which he should journey towards heaven. On this road
man is threatened by many dangers both from within and from without''
(Ps 141[142]:4). ``And therefore as guardians are appointed for men who have to
pass by an unsafe road, so an angel guardian is assigned to each man as long
as he is a wayfarer'' (\latin{quandiu viator est}). At the end of the road
the guardianship ends: ``he no longer has a guardian angel; but in the
kingdom he will have an angel to reign with him, in hell a demon to punish
him'' (\latin{sed habebit in regno Angelum conregnantem, in Inferno Daemonem
punientem}).

The replies take the three exceptions in turn. Christ as man ``was guided
immediately by the Word of God,'' and in his soul was a comprehensor, so
that ``it was right that He should have not a guardian angel as superior to
Him, but a ministering angel as inferior to Him,'' as when ``angels came and
ministered to Him'' (ad~1; Matt 4:11). The \emph{Sentences} commentary
answers here the Dionysian text that Christ ``readily subjected Himself to
the dispositions of the Father and God, through Angels'' (\emph{CH} 4.4): he
was subject to them indirectly, since the angels instructed Joseph and his
mother about the child (\emph{In II Sent.}\ d.~11 q.~1 a.~3 ad~6; Matt 2). Suárez adds that the Blessed Virgin had a guardian angel, citing
Bernard, while Christ had not a guardian but many ministers (\emph{De
angelis} VI.17.21). Adam in innocence suffered no danger from within, where
all was ordered, but danger threatened ``from without on account of the
snares of the demons; as was proved by the event. For this reason he needed a
guardian angel'' (ad~2). The foreknown, the infidels, and even the
Antichrist are not deprived of the inner help of natural reason, ``so neither
are they deprived of that exterior help granted by God to the whole human
race---namely the guardianship of the angels.'' If it does not bring them to
merit eternal life, it protects them ``from certain evils which would hurt
both themselves and others. For even the demons are held off by the good
angels, lest they hurt as much as they would. In like manner Antichrist
will not do as much harm as he would wish'' (ad~3). The earlier commentary
gives the reason in a phrase: the Antichrist will have a guardian
\latin{quia lex communis propter unum mutari non debet}, because the common
law is not to be changed for one man, and his punishment will appear the
more just because the gifts provided for all human nature were not taken
from him (\emph{In II Sent.}\ d.~11 q.~1 a.~3 ad~5).

Against those who restricted the guardians to the predestined or the just,
Suárez determines, \latin{non solum justos, sed etiam
peccatores, neque solos fideles, sed etiam infideles, neque solos baptizatos,
sed etiam non baptizatos, habere Angelos custodes}, not only the just but
sinners, not only the faithful but unbelievers, not only the baptized but the
unbaptized have guardian angels. \latin{Haec est communis sententia
theologorum et Patrum}; Basil, Hilary, and Chrysostom, who speak of
believers, add no exclusive word (\emph{De angelis} VI.17.12--14). The
guardian's proper office ends with death, \latin{quia omnia pericula
cessant, et spes profectus in morte finitur}, since every danger ceases and
the hope of progress ends; yet toward just souls an office of charity
remains. The guardians lead the pure to heaven, and lead those who need
purgation to the place of purgatory, where they visit and console them, as
the angels carried Lazarus; Suárez calls it the common sense of the Church,
\latin{valde credibile, et rationi consentaneum} (VI.19.9; Luke 16:22).

\paragraph{The guardian is given from birth (a.~5).} The objections argue
that men begin to receive the inheritance of salvation (Heb 1:14) at
baptism; that the newborn cannot be instructed; and that the child in the
womb has a rational soul and yet receives no sacrament. The \emph{sed contra} is
Jerome: ``each soul has an angel appointed to guard it from its birth.''
Aquinas begins from Origen: ``as Origen observes (Tract. v, super Matt.)
there are two opinions on this matter. For some have held that the angel
guardian is appointed at the time of baptism, others, that he is appointed
at the time of birth. The latter opinion Jerome approves \dots\ and with
reason.'' The reason is the distinction of gifts. ``Those benefits which are
conferred by God on man as a Christian, begin with his baptism; such as
receiving the Eucharist, and the like. But those which are conferred by God
on man as a rational being, are bestowed on him at his birth, for then it is
that he receives that nature. Among the latter benefits we must count the
guardianship of angels.'' The ministry of the angels is efficacious to
salvation for the heirs of salvation, but it is not withdrawn from others,
and it is efficacious for them too, ``inasmuch as they ward off many evils''
(ad~1). Instruction is the last and principal effect of the guardianship, but
it has others that suit childhood, ``to ward off the demons, and to prevent
both bodily and spiritual harm'' (ad~2). Of the child in the womb Aquinas
says that it is not yet wholly separate from the mother but ``by reason of a
certain intimate tie, is still part of her: just as the fruit while hanging
on the tree is part of the tree. And therefore it can be said with some
degree of probability, that the angel who guards the mother guards the child
while in the womb''; at birth, when it is separate, its own guardian is
appointed (ad~3).

The commentary on the \emph{Sentences} had placed the beginning earlier.
There the guardianship is owed to all men \latin{ab infusione animae
rationalis per quam ad finem salutis ordinantur \dots\ usque ad mortem},
from the infusion of the rational soul by which they are ordered to the end
of salvation, until death (\emph{In II Sent.}\ d.~11 q.~1 a.~3). Children in
the womb receive no sacraments because they are not subject to the acts of
ministers, but they are subject to the works of God and of the angels,
\latin{et ideo eis ab infusione animae rationalis, custos Angelus
assignatur}, and so a guardian angel is assigned to them from the infusion
of the rational soul, who holds back the demon's power and the many
hindrances by which the child's constitution might be made more prone to sin
or its life extinguished (ad~3). Honorius had said the same, that each soul
is committed to an angel \latin{dum in corpus mittitur}. Suárez sets the two
readings of Jerome's \latin{nativitas} side by side: birth from the womb, as
Aquinas in the \emph{Summa}, or the birth in the womb, conception, as the
\emph{Elucidarium}, Aquinas in the \emph{Sentences}, Bonaventure, and Gabriel
Biel. He argues for the second: from that moment man is a wayfarer and a
distinct person in the order of salvation, capable of original sin and, as
John the Baptist shows, of grace; and the mother, disobeying her own angel,
might attempt something harmful to the child that the child's own guardian
could prevent.
\latin{Denique haec providentia liberalior est, ideoque huic parti libentius
adhaereo}: finally, this providence is the more liberal, and so I adhere the
more willingly to this side (\emph{De angelis} VI.17.18). The gift belonged
to men under the law of nature and of Moses as under the law of grace, for
Jacob and David confess their angels, and the child Christ set among the
disciples did not yet belong to the law of grace (VI.17.19). The patristic question of birth or
baptism and its reception are set out in \S\ref{sec:dq-guardians}.

\paragraph{The guardian never wholly forsakes his charge (a.~6).} The
objections bring the angels' words over Babylon, ``let us forsake her'' (Jer
51:9), and the hedge taken from the vineyard, which the Gloss reads of ``the
guardianship of the angels'' (Isa 5:5); God himself forsakes man (Ps
21[22]:2), and the angel the more; and Damascene says that ``when the angels are here with us, they
are not in heaven,'' but sometimes they are in heaven (\emph{De fide orth.}\
II.3). The \emph{sed contra} argues from the enemy: ``The demons are ever assailing
us, according to 1 Peter 5:8: `Your adversary the devil, as a roaring lion,
goeth about, seeking whom he may devour.' Much more therefore do the good
angels ever guard us.'' Aquinas answers that the guardianship is ``an effect
of Divine providence in regard to man,'' and nothing is entirely withdrawn
from providence, since whatever shares in being is subject to it. God is said
to forsake a man only in permitting him to suffer some defect of punishment
or of fault. ``In like manner it must be said that the angel guardian never
forsakes a man entirely, but sometimes he leaves him in some particular, for
instance by not preventing him from being subject to some trouble, or even
from falling into sin, according to the ordering of Divine judgments''
(\latin{Angelus custos nunquam totaliter dimittit hominem, sed ad aliquid
interdum eum dimittit}). So Babylon and the house of Israel are said to be
forsaken by their angels, who did not keep them from tribulation. The angel
may leave a man as to place, but not as to the effect of his guarding, ``for
even when he is in heaven he knows what is happening to man; nor does he
need time for his local motion, for he can be with man in an instant''
(ad~3).

The earlier commentary adds three touches. The cry over Babylon may be
understood as \latin{vox Angelorum discedentium ab homine peccatore in hora
mortis; quia tunc primo desperatur de ejus salute}, the voice of the angels
departing from the sinner at the hour of death, when his salvation is first
despaired of (\emph{In II Sent.}\ d.~11 q.~1 a.~4 ad~1). A strong
impression remains in the one who receives it after the agent has gone, as
in violent motion, \latin{ut patet quando aliquis semel devote orat, ad
plures dies remanet inde devotior}, as when a man prays once devoutly and
remains the more devout for many days; so the angel, though not always
present, can always guard (ad~5). No one on the way is so obstinate that he
cannot be converted, and so the angel never gives him up (ad~6). Basil's image of the smoke
and the bees, and the Fathers who say that sinners drive their angels off,
Suárez explains of the actual exercise of the guardianship, which the angels
at times remit because of the sinner's incapacity or indisposition, not of
the care itself (\emph{De angelis} VI.17.15--16). He notes that the promise
of Exodus, ``he will not forgive when thou hast sinned'' (Exod 23:21), speaks
of leaving sin unpunished, not of deserting the sinner (VI.17.15). So
Bernard's guardian neither slumbers nor sleeps, though at times he seems to
look away (\emph{In Ps.\ Qui habitat} 12.9), and Gregory's rule for Job holds
of him: while guarding he forsakes, and while forsaking he guards
(\emph{Moralia} III, \S6).

\paragraph{The angels do not grieve for the ills of those they guard (a.~7).}
The objections: ``The angels of peace shall weep bitterly'' (Isa 33:7); the
loss of his charge is against the guardian's will; the angels who rejoice
over the penitent (Luke 15:7) should grieve over the fallen; and on the
first-fruits of Num 18:12 ``the gloss of Origen
says: `The angels are brought to judgment as to whether men have fallen
through their negligence or through their own fault','' and it is reasonable
to grieve for the ills that bring one to judgment. Origen's own sentence asks
whether the angel failed in his tilling or human sloth answered unworthily to
it (\emph{Hom.\ in Num.}\ XI.4; \S\ref{sec:gd-fathers}). The \emph{sed
contra} is the city of God: ``Death shall be no more, nor mourning, nor
crying, nor sorrow'' (Apoc 21:4); the angels are perfectly happy, and
``where there is grief and sorrow, there is not perfect happiness.''
Aquinas determines, \latin{Angeli non dolent neque de peccatis, neque de
poenis hominum}. Sorrow is only of what is contrary to the will, as
Augustine defines it; in the \emph{City of God} the pain of the soul called
sadness is ``a shrinking from those things which have happened to us in
spite of ourselves'' (\emph{De civ.\ Dei} XIV.15; cf.\ XIV.6). But nothing
happens in the world against the will of the angels and the blessed,
``because their will cleaves entirely to the ordering of Divine justice;
while nothing happens in the world save what is effected or permitted by
Divine justice.'' Aristotle's sailor who throws his cargo into the sea does
not will it absolutely and universally, but wills it in the danger to his
life; so ``universally and absolutely speaking the angels do not will sin
and the pains inflicted on its account: but they do will the fulfilment of
the ordering of Divine justice in this matter.'' Isaias's ``angels of
peace'' are literally the messengers of Ezechias who wept at the words of
Rabsaces (Isa 37:2), allegorically the apostles and preachers who weep for
men's sins; applied to the blessed angels the words are metaphorical and
signify ``that universally speaking the angels will the salvation of
mankind: for in this sense we attribute passions to God and the angels''
(ad~1). In repentance and in sin alike ``there is one reason for the angel's
joy, namely the fulfilment of the ordering of the Divine Providence'' (ad~3).
And ``the angels are brought into judgment for the sins of men, not as
guilty, but as witnesses to convict man of weakness'' (ad~4).

The \emph{Sentences} commentary refines the answer. The angels do not grieve,
but it is not for want of charity: \latin{compati non potest qui passibilis
non est; et ideo ex impassibilitate Angelorum hoc accidit quod condolere non
possunt, non ex caritatis defectu}, one who cannot suffer cannot suffer with
another (\emph{In II Sent.}\ d.~11 q.~1 a.~5 ad~2). Their joy is not lessened
by a man's loss, for the number of the elect cannot be lessened, \latin{Vel
potest dici, quod ipsi semper gaudent de suis bonis operibus quae custodiendo
egerunt, licet ille qui custoditus est non salvetur}, or because they always
rejoice in the good works they did in guarding, even if the one they guarded
is not saved (ad~3). They will his salvation with the antecedent will, as God
does, and with the consequent will his condemnation if he merits it (ad~4).
Bernard's angels who descend \latin{per compassionem tui} (\emph{In Ps.\ Qui
habitat} 11.6) show the same thing from the side of love: their compassion is
the mercy of their descent, and in descending they lose nothing of the
vision of glory. Suárez reads the Fathers who speak of angels grieving,
Chrysostom and Andrew of Caesarea among them, as the Scriptures ascribe
affections to God (\emph{De angelis} VI.19.12).

\paragraph{Strife among the angels (a.~8).} The objections: God ``maketh
peace in his high places'' (Job 25:2); perfect charity excludes strife; and
a good angel would favour injustice. The
\emph{sed contra} is Gabriel's word to Daniel: ``The prince of the kingdom of
the Persians resisted me one and twenty days,'' and ``this prince of the
Persians was the angel deputed to the guardianship of the kingdom of the
Persians. Therefore one good angel resists the others; and thus there is
strife among them'' (Dan 10:13). Aquinas sets two readings side by side.
Jerome, as Aquinas reports him, takes the prince for ``the angel who opposed
the setting free of the people of Israel,'' a resistance that ``may have been
caused by some prince of the demons having led the Jewish captives in Persia
into sin''; ``But according to Gregory (Moral. xvii), the prince of the
kingdom of Persia was a good angel appointed to the guardianship of that
kingdom.'' The divine judgments on kingdoms and men are carried out by the
angels; where contrary merits meet, the angels cannot know what the divine
wisdom has ordained ``unless God reveal it to them: and so they need to
consult Divine wisdom thereupon.'' ``Wherefore forasmuch as they consult the
Divine will concerning various contrary and opposing merits, they are said
to resist one another: not that their wills are in opposition, since they
are all of one mind as to the fulfilment of the Divine decree; but that the
things about which they seek knowledge are in opposition.'' The objections
are answered by this.

Jerome at his own place reads the prince as \latin{angelus cui Persis credita
est}, the angel to whom Persia was entrusted, who resisted \latin{faciens pro
credita sibi provincia}, acting for the province entrusted to him (\emph{In
Dan.}, on 10:13; PL~25; \S\ref{sec:fl-jerome}). In Gregory's exposition the
angels over the captives in Persia and over those left in Judaea bring the
claims of the peoples before ``the Court Above,'' and ``the one identical
victory of all of whom, however, is the Supreme Will of their Maker above
them'' (\emph{Moralia} XVII.12.17; \S\ref{sec:fa-gregory}). The
\emph{Sentences} commentary had weighed the two readings more sharply. There Jerome is
reported as holding the prince of Persia an evil angel set over the
Persians, \latin{sicut enim unicuique homini datur ad exercitium unus bonus
et alter malus Angelus; ita singulis gentibus duo spiritus praeponuntur, unus
bonus, et alter malus}: as each man has a good angel and an evil one for his
exercise, so two spirits are set over each nation, one good and one evil.
Aquinas objects that an evil angel could not have withstood a good
one so many days unless he had a just cause, and a just cause would be
better served by a good angel. \latin{Et ideo dicendum est cum Gregorio,
quod uterque bonus Angelus fuit: quod etiam magis litterae consonat}:
therefore we must say with Gregory that both were good angels, which agrees
better with the letter too, since it names in one series the prince of the
Persians and Michael the prince of the Jews. He adds a reason of his own
to the merits in conflict: Daniel's prayer deserved the people's release,
\latin{sed peccata populi, et etiam utilitas quam Judaei faciebant in regno
Persarum, dum per eos Dei notitia diffundebatur, erant in contrarium}; but
the sins of the people, and also the good the Jews were doing in the Persian
kingdom by spreading the knowledge of God, weighed on the other side. The
concord of the angels lies in their receiving the divine illumination, since
all will what they perceive God to will; and difference of opinion, as
Aristotle says, does not destroy friendship, but only difference of will
(\emph{In II Sent.}\ d.~11 q.~2 a.~5 and ad~1). Serenus in Cassian holds that
``this prince of the kingdom of the Persians was a hostile power, which
favoured the nation of the Persians an enemy of God's people'' (\emph{Conl.}\
VIII.13; \S\ref{sec:fl-cassian}). That the nations have their angels the
Fathers teach with one voice; for Damascene the angels ``are the guardians of
the divisions of the earth: they are set over nations and regions, allotted
to them by their Creator: they govern all our affairs and bring us succour''
(\emph{De fide orth.}\ II.3), and for Dionysius they lead their nations to
the one Head of all (\emph{CH} 9.4; \S\ref{sec:ch9}).

\angeldossier{\ST{I q.~113 aa.~1--8}; parallels \emph{In II Sent.}\ d.~11
q.~1 aa.~1--5 and q.~2 a.~5; \emph{SCG} III.80; \emph{De veritate} q.~8
a.~11.}
 {Church teaching: human life from infancy to death is surrounded by the
 angels' watchful care and intercession, and beside each believer stands an
 angel as protector and shepherd (CCC~336; Ps 90[91]:11; Matt 18:10; Heb
 1:14) (aa.~1--2); the Church venerates the guardian angels with a feast and
 commends prayer to them (John Paul~II, 6 August 1986; \emph{Directory}
 215--216). Common teaching: every man, believer or not, has his own
 guardian for the whole of his wayfaring (a.~4; Jerome; Chrysostom;
 Suárez, \emph{De angelis} VI.17.14); the guardian is given from the
 beginning of life, not from baptism (a.~5; Jerome); he never wholly
 forsakes his charge in this life (a.~6); the guarding of single men belongs
 to the lowest order and that of peoples to the Principalities or
 Archangels (a.~3; Dionysius, \emph{CH}~9; Gregory, \emph{Hom.\ in Ev.}\ 34);
 the blessed angels suffer no sorrow at men's sins or pains (a.~7; Apoc
 21:4); angels preside over nations (a.~8; Deut 32:8 LXX; Dan 10). Thomist
 position: the greater angels probably guard those chosen for greater glory
 (a.~3 ad~1); Christ had ministering angels and no guardian (a.~4 ad~1); the
 ``resistance'' of Dan 10:13 is the consulting of the divine will on contrary
 merits (a.~8). Disputed: whether the guardian is given at birth, the mother's
 angel guarding the child in the womb (a.~5 ad~3), or at the infusion of the
 soul (\emph{In II Sent.}\ d.~11 q.~1 a.~3 ad~3; Honorius; Suárez VI.17.18);
 in the patristic tradition, at birth or at baptism (Origen, \emph{Comm.\ in
 Matt.}\ XIII.27); whether the prince of Persia was a good angel (Gregory;
 Aquinas) or an evil or hostile one (Jerome as reported in \emph{In II
 Sent.}\ d.~11 q.~2 a.~5; Cassian). Pious tradition: the guardians lead just
 souls to heaven, or to purgatory where they console them (Suárez VI.19.9;
 Luke 16:22).}
 {\emph{Sed contra}: Ps 90[91]:11 (aa.~1, 3), Jerome, \emph{In Matth.}\ III on
 18:10 (aa.~2, 4, 5), 1 Pet 5:8 (a.~6), Apoc 21:4 (a.~7), Dan 10:13 (a.~8).
 In the articles: Chrysostom, \emph{Hom.\ in Matt.}\ 59.4; Dionysius,
 \emph{CH} 4, 5, 9, 10; Gregory, \emph{Hom.\ in Ev.}\ 34, \emph{Moralia}
 XVII.12.17; Origen on Matthew and the Gloss from \emph{Hom.\ in Num.}\ XI.4;
 Damascene, \emph{De fide orth.}\ II.3; Augustine, \emph{De civ.\ Dei}
 XIV.15; Aristotle, \emph{Eth.}\ III.1. Read for the doctrine: Hermas,
 \emph{Mand.}\ VI.2; Clement, \emph{Strom.}\ VI.17; Origen, \emph{Comm.\ in
 Matt.}\ XIII.26--28, \emph{Hom.\ in Num.}\ XI.3--4, \emph{De princ.}\ II.10.7;
 Basil, \emph{Adv.\ Eun.}\ III.1, \emph{Hom.\ in Ps.}\ 33.5; Gregory of Nyssa,
 \emph{De vita Moysis} II; Chrysostom, \emph{Hom.\ in Act.}\ 26, \emph{Hom.\
 in Col.}\ 3; Hilary, \emph{Tract.\ in Ps.}\ 124.6, 134.17; Ambrose, \emph{De
 viduis} 9.55; Cassian, \emph{Conl.}\ VIII.13, 17; Isidore, \emph{Sent.}\
 I.10.20--21; Bernard, \emph{In Ps.\ Qui habitat} 11--13; Honorius,
 \emph{Elucidarium} II.28--29; Peter Lombard, \emph{Sent.}\ II d.~11 c.~1;
 Suárez, \emph{De angelis} VI.17--19; CCC~336; John Paul~II, 6 August 1986;
 \emph{Directory on Popular Piety} 215--217.}
 {Assignment of the guardian at baptism (a.~5 obj.~1); Origen's second
 exposition, that until regeneration men are under the angels of Satan and no
 holy angel attends them (\emph{Comm.\ in Matt.}\ XIII.28); Chrysostom's
 reading of Matt 18:10 of greater angels (a.~3 obj.~1, answered in ad~1); the
 Lombard's one angel for many men at once (\emph{Sent.}\ II d.~11 c.~1),
 answered in a.~2; restrictions of the guardians to the predestined, to the
 just, or to believers, rejected by Suárez (VI.17.12--14) and by a.~4 ad~3;
 the \emph{Sentences} commentary's guardian from the infusion of the soul,
 against a.~5 ad~3; Jerome's evil or Cassian's hostile prince of Persia,
 against Gregory (a.~8); Calvin, as Suárez reports him, doubting a proper
 angel for each man (VI.17.8--9).}

\sectionguard
\subsection{The Guardian Angel in the Witnesses}\label{sec:gd-witnesses}

From the psalm to the Catechism the witnesses agree on the angel of each
soul; they divide on the moment of its gift, on the rank of the angel who
guards, and on the prince of Persia.

\begin{fourstable}{Witness}{Locus}{Teaching on the guardian angel}{Grade}
Scripture & Ps 90[91]:11--12; Matt 18:10; Tob 5--12; Heb 1:14 & God gives his angels charge to keep man in all his ways; the little ones have angels who see the Father's face & Church teaching (CCC~336)\\
Scripture & Deut 32:8 (LXX); Dan 10:13, 20--21; 12:1 & Angels preside over nations; Michael over Israel & Common teaching\\
Hermas & \emph{Mand.}\ VI.2 & Two angels are with each man, known by the thoughts they raise & Common teaching\\
Origen & \emph{Comm.\ in Matt.}\ XIII.26--28 & The little ones have angels as nurses and tutors; from birth or from regeneration & Disputed (the moment)\\
Origen & \emph{Hom.\ in Num.}\ XI.3--4 & The angels till human hearts and bring their charges to the judgment & Common teaching (a.~7 ad~4)\\
Basil & \emph{Adv.\ Eun.}\ III.1 & Beside each believer an angel as tutor and shepherd; the angel of a nation is greater & Church teaching (CCC~336)\\
Basil & \emph{Hom.\ in Ps.}\ 33.5 & Sin drives the guardian off as smoke drives bees; one angel is a whole camp & Common teaching\\
Gregory of Nyssa & \emph{De vita Moysis} II & After the fall an angel was set beside each life as ally, a demon opposite; the angel is brother to the mind & Common teaching\\
Chrysostom & \emph{Hom.\ in Matt.}\ 59.4; \emph{Act.}\ 26; \emph{Col.}\ 3 & Each man has an angel; now the angels are by believers; the angels of Matt 18:10 are greater & Common teaching\\
Hilary & \emph{Tract.\ in Ps.}\ 124.6; 134.17 & Angels given to guard a weak nature; the Lord's guard is better & Common teaching\\
Ambrose & \emph{De viduis} 9.55 & The angels given as guards are to be entreated for us & Church teaching (CCC~335)\\
Jerome & \emph{In Matth.}\ III (18:10) & Each soul has an angel from its birth & Common teaching (aa.~2, 4--5 s.c.)\\
Cassian & \emph{Conl.}\ VIII.17 & A good and a bad angel cling to each man & Common teaching\\
Gregory the Great & \emph{Hom.\ in Ev.}\ 34.8, 10; \emph{Mor.}\ XVII.12.17 & Angels announce the least things; the princes of nations judge their merits & Common teaching\\
Honorius & \emph{Elucidarium} II.28--29 & Each soul is committed to an angel when sent into the body; the angels come when invoked & Disputed (the moment); Common teaching\\
Peter Lombard & \emph{Sent.}\ II d.~11 c.~1 & Each has a good angel for guard, an evil for exercise; one angel may guard many & Common teaching; answered in a.~2\\
Bernard & \emph{In Ps.\ Qui habitat} 11--13 & Reverence for the angel's presence, devotion for his goodwill, confidence in his guard & Common teaching; Church teaching (\emph{Directory} 216)\\
Aquinas & \ST{I q.~113} & Guardianship executes providence toward each immortal person, from birth, never wholly withdrawn & Common teaching; Thomist position\\
Aquinas & \emph{In II Sent.}\ d.~11 q.~1 a.~3 & The guardian is given at the infusion of the rational soul & Disputed\\
Suárez & \emph{De angelis} VI.17--19 & Guardianship certain of faith; from conception; guardians lead just souls to heaven or purgatory & Common teaching; Pious tradition\\
Catechism; John Paul~II; \emph{Directory} & CCC~336; audience of 6 Aug.\ 1986; nn.~215--217 & Human life is surrounded by the angels' care; the Church keeps their feast and prays the \latin{Angele Dei} & Church teaching\\
\end{fourstable}
